% Part: first-order-logic
% Chapter: syntax-and-semantics
% Section: satisfaction

\documentclass[../../../include/open-logic-section]{subfiles}

\begin{document}

\olfileid[pt]{fol}{syn}{sat}

\olsection{Satisfação de \article{formula} \printtoken{S}{formula}
  em \article{structure} \printtoken{S}{structure}}

\begin{explain}
A noção básica que relaciona expressões como termos e
!!{formula}s, por um lado, e !!{structure}s, por outro, são
as de \emph{!!{value}} de um termo e \emph{satisfação} de
!!a{formula}.  Informalmente, o !!{value} de um termo é !!a{element} de
!!a{structure}---se o termo é apenas uma constante, seu !!{value} é o
objeto atribuído à constante pela !!{structure}, e se ele é
construído usando !!{function}s, o !!{value} é calculado a partir dos
!!{value}s das constantes e das funções atribuídas às funções no
termo.  !!^a{formula} é \emph{satisfeita} em !!a{structure} se a
interpretação dada aos predicados torna a !!{formula} verdadeira no
domínio da !!{structure}. Essa noção de satisfação é
especificada indutivamente: a especificação da !!{structure} afirma
diretamente quando !!{formula}s atômicas são satisfeitas, e definimos quando uma
!!{formula} complexa é satisfeita dependendo do conectivo principal ou
quantificador e de se as !!{subformula}s imediatas são ou não
satisfeitas.

O caso dos quantificadores aqui é um pouco delicado, pois a
!!{subformula} imediata de uma !!{formula} quantificada tem uma
!!{variable} livre, e !!{structure}s não especificam os !!{value}s das
!!{variable}s.  Para lidar com essa dificuldade, também
introduzimos \emph{atribuições de variáveis} e definimos a satisfação não em
relação a !!a{structure} sozinha, mas em relação a !!a{structure}
mais uma atribuição de variáveis.
\end{explain}

\begin{defn}[Atribuição de Variáveis]
Uma \emph{atribuição de variáveis}~$s$ para !!a{structure}~$\Struct{M}$ é uma
função que mapeia cada !!{variable} para !!a{element} de~$\Domain M$,
isto é, $s\colon \Var \to \Domain M$.
\end{defn}

\begin{explain}
!!^a{structure} atribui !!a{value} a cada !!{constant}, e uma
atribuição de variáveis a cada variável.  Mas queremos usar termos construídos
a partir deles para também nomear !!{element}s do !!{domain}.  Para isso,
definimos o !!{value} dos termos indutivamente. Para !!{constant}s e
variáveis, o valor é exatamente como a !!{structure} ou a atribuição de
variáveis o especifica; para termos mais complexos, ele é calculado
recursivamente usando as funções que a !!{structure} atribui aos
!!{function}s.
\end{explain}

\begin{defn}[!!^{value} de Termos]
Se $t$ é um termo da linguagem~$\Lang L$, $\Struct M$ é
!!a{structure} para~$\Lang L$, e $s$ é uma atribuição de variáveis
para~$\Struct M$, o \emph{!!{value}}~$\Value{t}{M}[s]$ é definido como
segue:
\begin{enumerate}
\item \indcase{t}{c}{$\Value{\indfrm}{M}[s] = \Assign{\indcomplex}{M}$.}
\item \indcase{t}{x}{$\Value{\indfrm}{M}[s] = s(\indcomplex)$.}
\item \indcase{t}{\Atom{f}{t_1, \ldots, t_n}}{
\[
\Value{\indfrm}{M}[s] = \Assign{f}{M}(\Value{t_1}{M}[s], \ldots,
\Value{t_n}{M}[s]).
\]}
\end{enumerate}
\end{defn}

\begin{defn}[$x$-Variante]
Se $s$ é uma atribuição de variáveis para !!a{structure}~$\Struct M$, então qualquer
atribuição de variáveis~$s'$ para~$\Struct M$ que difere de~$s$ no máximo
no que atribui a~$x$ é chamada de \emph{$x$-variante} de~$s$.  Se
$s'$ é uma $x$-variante de~$s$ escrevemos $\varAssign{s'}{s}{x}$.
\end{defn}

\begin{explain}
Note que uma $x$-variante de uma atribuição~$s$ não \emph{precisa}
atribuir algo diferente a~$x$.  De fato, toda atribuição conta
como uma $x$-variante de si mesma.
\end{explain}

\begin{defn}
  Se $s$ é uma atribuição de variáveis para !!a{structure}~$\Struct M$
  e $m \in \Domain{M}$, então a atribuição~$\Subst{s}{m}{x}$ é a
  atribuição de variáveis definida por
  \[\Subst{s}{m}{x}(y) = \begin{cases}
    m & \text{se } y \ident x\\
    s(y) & \text{caso contrário}.
  \end{cases}\]
\end{defn}

Em outras palavras, $\Subst{s}{m}{x}$ é a $x$-variante particular de~$s$
que atribui o !!{element}~$m$ do domínio a~$x$, e atribui as mesmas
coisas às !!{variable}s diferentes de~$x$ que $s$ atribui.

\begin{defn}[Satisfação]
\ollabel{defn:satisfaction}
A satisfação de !!a{formula}~$!A$ em !!a{structure}~$\Struct M$
relativamente a uma atribuição de variáveis~$s$, em símbolos:
$\Sat{M}{!A}[s]$, é definida recursivamente como segue. (Escrevemos
$\Sat/{M}{!A}[s]$ para significar ``não $\Sat{M}{!A}[s]$''.)
\begin{enumerate}
\tagitem{prvFalse}{%
  \indcase{!A}{\lfalse}{$\Sat/{M}{\indfrm}[s]$.}}{}

\tagitem{prvTrue}{%
  \indcase{!A}{\ltrue}{$\Sat{M}{\indfrm}[s]$.}}{}

\item \indcase{!A}{\Atom{R}{t_1, \dots, t_n}}{$\Sat{M}{\indfrm}[s]$
  sse $\langle \Value{t_1}{M}[s], \dots, \Value{t_n}{M}[s] \rangle \in
  \Assign{R}{M}$.}

\item \indcase{!A}{\eq[t_1][t_2]}{$\Sat{M}{\indfrm}[s]$ sse
  $\Value{t_1}{M}[s] = \Value{t_2}{M}[s]$.}

\tagitem{prvNot}{%
  \indcase{!A}{\lnot !B}{$\Sat{M}{\indfrm}[s]$ sse
    $\Sat/{M}{!B}[s]$.}}{}

\tagitem{prvAnd}{%
  \indcase{!A}{(!B \land !C)}{$\Sat{M}{\indfrm}[s]$ sse $\Sat{M}{!B}[s]$
    e $\Sat{M}{!C}[s]$.}}{}

\tagitem{prvOr}{%
  \indcase{!A}{(!B \lor !C)}{$\Sat{M}{\indfrm}[s]$ sse
    $\Sat{M}{!B}[s]$ ou $\Sat{M}{!C}[s]$ (ou ambos).}}{}

\tagitem{prvIf}{%
  \indcase{!A}{(!B \lif !C)}{$\Sat{M}{\indfrm}[s]$ sse $\Sat/{M}{!B}[s]$
    ou $\Sat{M}{!C}[s]$ (ou ambos).}}{}

\tagitem{prvIff}{%
  \indcase{!A}{(!B \liff !C)}{$\Sat{M}{\indfrm}[s]$ sse ou ambos
    $\Sat{M}{!B}[s]$ e $\Sat{M}{!C}[s]$, ou nem $\Sat{M}{!B}[s]$
    nem $\Sat{M}{!C}[s]$.}}{}

\tagitem{prvAll}{%
  \indcase{!A}{\lforall[x][!B]}{$\Sat{M}{\indfrm}[s]$ sse para todo
    !!{element}~$m \in \Domain M$, $\Sat{M}{!B}[\Subst{s}{m}{x}]$.}}{}

\tagitem{prvEx}{%
  \indcase{!A}{\lexists[x][!B]}{$\Sat{M}{\indfrm}[s]$ sse para pelo menos
  um !!{element}~$m \in \Domain M$, $\Sat{M}{!B}[\Subst{s}{m}{x}]$.}}{}
\end{enumerate}
\end{defn}

\begin{explain}
As atribuições de variáveis são importantes nas últimas
\iftag{notprvEx,notprvAll}{cláusula}{duas cláusulas}.\iftag{prvAll}{ Não
podemos definir a satisfação de $\lforall[x][!B(x)]$ por ``para todo $m \in
\Domain{M}$, $\Sat{M}{!B(m)}$''.}{}\iftag{prvEx}{ Não podemos definir a
satisfação de $\lexists[x][!B(x)]$ por ``para pelo menos um $m \in
\Domain{M}$, $\Sat{M}{!B(m)}$''.}{} A razão é que, se $m \in
\Domain M$, ele não é um símbolo da linguagem, e portanto $!B(m)$~não é
!!a{formula} (isto é, $\Subst{!B}{m}{x}$ é indefinido).  Também não
podemos supor que temos !!{constant}s ou termos disponíveis que nomeiem
cada !!{element} de~$\Struct{M}$, já que não há nada na
definição de !!{structure}s que o exija.  Na linguagem padrão,
o conjunto de !!{constant}s é !!{denumerable}, então se
$\Domain{M}$ não é !!{enumerable} não há sequer !!{constant}s
suficientes para nomear cada objeto.

Resolvemos esse problema introduzindo atribuições de variáveis, que
nos permitem ligar variáveis diretamente a !!{element}s do domínio.
Então, em vez de dizer que, por exemplo,
\iftag{prvEx}{$\lexists[x][!B(x)]$}{$\lforall[x][!B(x)]$} é satisfeita
em~$\Struct M$ sse \iftag{prvEx}{para pelo menos um $m \in
\Domain{M}$}{para todo $m \in \Domain{M}$, $\Sat{M}{!B(m)}$}, dizemos que ela
é satisfeita em~$\Struct M$ \emph{relativamente a}~$s$ sse $!B(x)$ é
satisfeita relativamente a~$\Subst{s}{m}{x}$ \iftag{prvEx}{para pelo
menos um}{para todo} $m \in \Domain M$.
\end{explain}

\begin{ex}
Seja $\Lang{L} = \{a, b, f, R\}$ onde $a$ e $b$ são !!{constant}s,
$f$~é um !!{function} de dois lugares, e $R$~é um !!{predicate} de dois lugares.
Considere a !!{structure}~$\Struct{M}$ definida por:
\begin{enumerate}
\item $\Domain M = \{1, 2, 3, 4\}$
\item $\Assign{a}{M} = 1$
\item $\Assign{b}{M} = 2$
\item $\Assign{f}{M}(x, y) = x+y$ se $x+y \le 3$ e $= 3$ caso contrário.
\item $\Assign{R}{M} = \{\tuple{1, 1}, \tuple{1, 2}, \tuple{2, 3}, \tuple{2, 4}\}$
\end{enumerate}
A função $s(x) = 1$ que atribui $1 \in \Domain{M}$ a cada
!!{variable} é uma atribuição de variáveis para~$\Struct{M}$.

Então
\begin{align*}
\Value{f(a,b)}{M}[s] & = \Assign{f}{M}(\Value{a}{M}[s], \Value{b}{M}[s]).
\intertext{Como $a$ e $b$ são !!{constant}s, $\Value{a}{M}[s]
  = \Assign{a}{M} = 1$ e $\Value{b}{M}[s] = \Assign{b}{M} = 2$. Então}
\Value{f(a,b)}{M}[s] & = \Assign{f}{M}(1, 2) = 1+2 = 3.
\intertext{Para calcular o valor de $f(f(a,b),a)$ temos que considerar}
\Value{f(f(a,b),a)}{M}[s] & = \Assign{f}{M}(\Value{f(a, b)}{M}[s],
\Value{a}{M}[s]) = \Assign{f}{M}(3, 1) = 3,
\intertext{já que $3+1 > 3$. Como $s(x) = 1$ e $\Value{x}{M}[s] =
  s(x)$, também temos}
\Value{f(f(a,b),x)}{M}[s] & = \Assign{f}{M}(\Value{f(a, b)}{M}[s],
\Value{x}{M}[s]) = \Assign{f}{M}(3, 1) = 3,
\end{align*}

Uma !!{formula} atômica~$R(t_1, t_2)$ é satisfeita se a tupla de
valores de seus argumentos, isto é, $\tuple{\Value{t_1}{M}[s],
  \Value{t_2}{M}[s]}$, é !!a{element} de~$\Assign{R}{M}$. Assim, por exemplo,
temos $\Sat{M}{R(b,f(a,b))}[s]$ já que $\tuple{\Value{b}{M},
  \Value{f(a,b)}{M}} = \tuple{2, 3} \in \Assign{R}{M}$, mas
$\Sat/{M}{R(x, f(a,b))}[s]$ já que $\tuple{1, 3} \notin \Assign{R}{M}[s]$.

Para determinar se uma fórmula não atômica~$!A$ é satisfeita, aplicam-se as
cláusulas da definição indutiva que se aplicam ao conectivo
principal. Por exemplo, o conectivo principal em $R(a, a) \lif (R(b,
x) \lor R(x, b))$ é o~$\lif$, e
\begin{align*}
  & \Sat{M}{R(a, a) \lif (R(b, x) \lor R(x, b))}[s] \text{ sse }\\
  & \qquad
  \Sat/{M}{R(a,a)}[s] \text{ ou } \Sat{M}{R(b, x) \lor R(x, b)}[s]
  \intertext{Como $\Sat{M}{R(a,a)}[s]$ (porque $\tuple{1,1} \in
    \Assign{R}{M}$) ainda não podemos determinar a resposta e devemos primeiro
    descobrir se $\Sat{M}{R(b, x) \lor R(x, b)}[s]$:}
  & \Sat{M}{R(b, x) \lor R(x, b)}[s] \text{ sse }\\
  & \qquad \Sat{M}{R(b,x)}[s] \text{ ou } \Sat{M}{R(x, b)}[s]
  \intertext{E este é o caso, já que $\Sat{M}{R(x, b)}[s]$
    (porque $\tuple{1,2} \in \Assign{R}{M}$).}
\end{align*}

Lembre-se de que uma $x$-variante de~$s$ é uma atribuição de variáveis que
difere de $s$ no máximo no que atribui a~$x$. Para cada
!!{element} de~$\Domain{M}$, há uma $x$-variante de~$s$:
\begin{align*}
  s_1 & = \Subst{s}{1}{x}, &
  s_2 & = \Subst{s}{2}{x},\\
  s_3 & = \Subst{s}{3}{x}, &
  s_4 & = \Subst{s}{4}{x}.
\end{align*}
Assim, por exemplo, $s_2(x) = 2$ e $s_2(y) = s(y) = 1$ para todas as variáveis~$y$
diferentes de~$x$. Estas são todas as $x$-variantes de~$s$ para a
estrutura~$\Struct{M}$, já que $\Domain{M} = \{1, 2, 3, 4\}$. Note, em
particular, que $s_1 = s$ ($s$~é sempre uma $x$-variante de si mesma).

\iftag{prvEx}{Para determinar se uma !!{formula} quantificada
  existencialmente~$\lexists[x][!A(x)]$ é satisfeita, temos que determinar
  se $\Sat{M}{!A(x)}[\Subst{s}{m}{x}]$ para pelo menos um $m \in \Domain
  M$. Assim,
  \[
  \Sat{M}{\lexists[x][(R(b,x) \lor R(x,b))]}[s],
  \]
  já que $\Sat{M}{R(b,x) \lor R(x, b)}[\Subst{s}{1}{x}]$
  ($\Subst{s}{3}{x}$ também serviria).  Mas,
  \[
  \Sat/{M}{\lexists[x][(R(b,x) \land R(x,b))]}[s]
  \]
  já que, qualquer que seja o $m \in \Domain{M}$ que escolhamos, $\Sat/{M}{R(b,x) \land R(x,b)}[\Subst{s}{m}{x}]$.}{}

\iftag{prvAll}{Para determinar se uma !!{formula} quantificada
  universalmente~$\lforall[x][!A(x)]$ é satisfeita, temos que determinar
  se $\Sat{M}{!A(x)}[\Subst{s}{m}{x}]$ para todo $m \in \Domain M$. Assim,
  \[
  \Sat{M}{\lforall[x][(R(x,a) \lif R(a,x))]}[s],
  \]
  já que $\Sat{M}{R(x,a) \lif R(a,x)}[\Subst{s}{m}{x}]$ para todo $m \in
  \Domain M$. Para $m = 1$, temos $\Sat{M}{R(a,x)}[\Subst{s}{1}{x}]$
  então o consequente é verdadeiro; para $m = 2$, $3$, e~$4$, temos
  $\Sat/{M}{R(x,a)}[\Subst{s}{m}{x}]$, então o antecedente é falso.
  Mas,
  \[
  \Sat/{M}{\lforall[x][(R(a,x) \lif R(x,a))]}[s]
  \]
  já que $\Sat/{M}{R(a,x) \lif R(x,a)}[\Subst{s}{2}{x}]$ (porque
  $\Sat{M}{R(a, x)}[\Subst{s}{2}{x}]$ e $\Sat/{M}{R(x,
  a)}[\Subst{s}{2}{x}]$).}{}

\iftag{defEx}{Para determinar se uma !!{formula} quantificada
  existencialmente~$\lexists[x][!A(x)]$ é satisfeita, temos que determinar
  se $\Sat{M}{\lnot \lforall[x][\lnot !A(x)]}[s]$. Por exemplo, temos
  \[
  \Sat{M}{\lexists[x][(R(b,x) \lor R(x,b))]}[s].
  \]
  Primeiro, $\Sat{M}{R(b,x) \lor R(x, b)}[\Subst{s}{1}{x}]$
  ($\Subst{s}{3}{x}$ também serviria).  Assim,
  $\Sat/{M}{\lnot(R(b,x) \lor R(x,b))}[\Subst{s}{1}{x}]$, logo
  $\Sat/{M}{\lforall[x][\lnot((R(b,x) \lor R(x,b))]}[s]$, e
  portanto $\Sat{M}{\lnot\lforall[x][\lnot((R(b,x) \lor
  R(x,b))]}[s]$. Por outro lado,
  \[
  \Sat/{M}{\lexists[x][(R(b,x) \land R(x,b))],}[s].
  \]
  Isso porque $\Sat{M}{\lforall[x][\lnot(R(b,x) \land
  R(x,b))]}[s]$, já que para nenhum $m \in \Domain M$, $\Sat{M}{R(b,x) \land
  R(x,b)}[\Subst{s}{m}{x}]$. Como você provavelmente pode adivinhar a partir destes
  exemplos, $\Sat{M}{\lexists[x][!A(x)]}[s]$ sse
  $\Sat{M}{!A(x)}[\Subst{s}{m}{x}]$ para pelo menos um $m \in \Domain
  M$.}{}

\iftag{defAll}{Para determinar se uma !!{formula} quantificada
  universalmente~$\lforall[x][!A(x)]$ é satisfeita, temos que determinar
  se $\Sat{M}{\lnot\lexists[x][\lnot !A(x)]}[s]$. Por exemplo,
  \[
  \Sat{M}{\lforall[x][(R(x,a) \lif R(a,x))]}[s],
  \]
  Primeiro, $\Sat{M}{R(x,a) \lif R(a,x)}[\Subst{s}{m}{x}]$ para todo $m \in
  \Domain M$ ($\Sat{M}{R(a,x)}[\Subst{s}{1}{x}]$ e
  $\Sat/{M}{R(a,x)}[\Subst{s}{m}{x}]$ para $m = 2$, $3$, ou~$4$).
  Assim, não há $m \in \Domain M$ tal que
  $\Sat{M}{\lnot(R(x,a) \lif R(a,x))}[\Subst{s}{m}{x}]$ e portanto
  $\Sat/{M}{\lexists[x][\lnot(R(x,a) \lif R(a,x))]}[s]$. Portanto,
  $\Sat{M}{\lnot\lexists[x][\lnot(R(x,a) \lif R(a,x))]}[s]$. Por outro
  lado,
  \[
  \Sat/{M}{\lforall[x][(R(a,x) \lif R(x,a))]}[s],
  \]
  já que $\Sat/{M}{R(a,x) \lif R(x,a)}[\Subst{s}{2}{x}]$, e assim
  $\Sat{M}{\lexists[x][\lnot(R(a,x) \lif R(x,a))]}[s]$. Como você pode
  provavelmente adivinhar a partir destes exemplos, $\Sat{M}{\lforall[x][!A(x)]}[s]$
  sse $\Sat{M}{!A(x)}[\Subst{s}{m}{x}]$ para todo $m \in \Domain
  M$.}{}

Para um caso mais complicado, considere
\[
\lforall[x][(R(a,x) \lif \lexists[y][R(x,y)])].
\]
Como $\Sat/{M}{R(a,x)}[\Subst{s}{3}{x}]$ e
$\Sat/{M}{R(a,x)}[\Subst{s}{4}{x}]$, os casos interessantes em que
temos que nos preocupar com o consequente do condicional são apenas $m = 1$
e $ = 2$. $\Sat{M}{\lexists[y][R(x,y)]}[\Subst{s}{1}{x}]$
vale? Vale se há pelo menos um $n \in \Domain M$ tal que
$\Sat{M}{R(x,y)}[\Subst{\Subst{s}{1}{x}}{n}{y}]$. De fato, se tomarmos
$n = 1$, temos $\Subst{\Subst{s}{1}{x}}{n}{y} = \Subst{s}{1}{y} =
s$. Como $s(x) = 1$, $s(y) = 1$, e $\tuple{1,1} \in \Assign{R}{M}$,
a resposta é sim.

Para determinar se $\Sat{M}{\lexists[y][R(x,y)]}[\Subst{s}{2}{x}]$,
temos que olhar para as atribuições de variáveis
$\Subst{\Subst{s}{2}{x}}{n}{y}$. Aqui, para $n = 1$, essa atribuição
é~$s_2 = \Subst{s}{2}{x}$, que não satisfaz $R(x,y)$ ($s_2(x) =
2$, $s_2(y) = 1$, e $\tuple{2,1}\notin \Assign{R}{M}$). No entanto,
considere $\Subst{\Subst{s}{2}{x}}{3}{y} = \Subst{s_2}{3}{y}$.
$\Sat{M}{R(x,y)}[\Subst{s_2}{3}{y}]$ já que $\tuple{2,3} \in
\Assign{R}{M}$, e assim $\Sat{M}{\lexists[y][R(x,y)]}[s_2]$.

Assim, para todo $n \in \Domain M$, ou
$\Sat/{M}{R(a,x)}[\Subst{s}{m}{x}]$ (se $m = 3$, $4$) ou
$\Sat{M}{\lexists[y][R(x,y)]}[\Subst{s}{m}{x}]$ (se $m = 1$, $2$), e assim
\[
\Sat{M}{\lforall[x][(R(a,x) \lif \lexists[y][R(x,y)])]}[s].
\]
Por outro lado,
\[
\Sat/{M}{\lexists[x][(R(a,x) \land \lforall[y][R(x,y)])]}[s].
\]
Temos $\Sat{M}{R(a,x)}[\Subst{s}{m}{x}]$ apenas para $m = 1$ e $m =
2$. Mas para ambos esses valores de~$m$, há por sua vez um $n \in
\Domain M$, a saber $n = 4$, tal que
$\Sat/{M}{R(x,y)}[\Subst{\Subst{s}{m}{x}}{n}{y}]$ e assim
$\Sat/{M}{\lforall[y][R(x,y)]}[\Subst{s}{m}{x}]$ para $m = 1$ e $m =
2$. Em suma, não há $m \in \Domain M$ tal que $\Sat{M}{R(a,x)
\land \lforall[y][R(x,y)]}[\Subst{s}{m}{x}]$.
\end{ex}

\iftag{defEx}{%
\begin{prop}\ollabel{prop:sat-ex}
$\Sat{M}{\lexists[x][!B(x)]}[s]$ sse há uma $x$-variante $s'$ de $s$
    tal que $\Sat{M}{!B(x)}[s']$.
\end{prop}

\begin{proof}
  Exercício.
\end{proof}
}{}

\tagprob{defEx}
\begin{prob}
  Prove \olref[fol][syn][sat]{prop:sat-ex}
\end{prob}
\tagendprob

\iftag{defAll}{%
\begin{prop}\ollabel{prop:sat-all}
$\Sat{M}{\lforall[x][!B(x)]}[s]$ sse para toda $x$-variante~$s'$ de $s$,
  $\Sat{M}{!B(x)}[s']$
\end{prop}

\begin{proof}
  Exercício.
\end{proof}
}{}

\tagprob{defAll}
\begin{prob}
  Prove \olref[fol][syn][sat]{prop:sat-all}
\end{prob}
\tagendprob

\begin{prob}
Seja $\Lang L = \{c, f, A\}$ com um !!{constant}, um !!{function}
de um lugar e um !!{predicate} de dois lugares, e seja a
!!{structure}~$\Struct{M}$ dada por
\begin{enumerate}
\item $\Domain M = \{1, 2, 3\}$
\item $\Assign{c}{M} = 3$
\item $\Assign{f}{M}(1) = 2, \Assign{f}{M}(2) = 3, \Assign{f}{M}(3) = 2$
\item $\Assign{A}{M} = \{\tuple{1, 2}, \tuple{2, 3}, \tuple{3, 3}\}$
\end{enumerate}
(a) Seja $s(v) = 1$ para todas as !!{variable}s~$v$.  Descubra se
\[
\Sat{M}{\lexists[x][(A(f(z), c) \lif \lforall[y][(A(y, x) \lor A(f(y),
      x))])]}[s]
\]
Explique por que sim ou por que não.

(b) Dê uma estrutura e uma atribuição de variáveis diferentes nas quais a
!!{formula} não seja satisfeita.
\end{prob}

\end{document}
